\documentclass[11pt]{article}
\usepackage[T1]{fontenc}
\usepackage{lmodern,amsmath,amssymb,amsthm,mathtools,microtype,booktabs}
\usepackage[margin=1in]{geometry}
\usepackage[hidelinks]{hyperref}
\usepackage{needspace}
\setlength{\emergencystretch}{2em}
\newtheorem{theorem}{Theorem}[section]
\newtheorem{lemma}[theorem]{Lemma}
\newtheorem{proposition}[theorem]{Proposition}
\newtheorem{corollary}[theorem]{Corollary}
\theoremstyle{remark}\newtheorem{remark}[theorem]{Remark}
\newcommand{\R}{\mathbb R}\newcommand{\Var}{\operatorname{Var}}
\newcommand{\Ptwo}{\mathcal P_2}\newcommand{\TV}{\operatorname{TV}}
\newcommand{\norm}[1]{\lVert#1\rVert}\newcommand{\ind}{\mathbf 1}
\title{Tail-sensitive convex-gradient interpolation\\and Brenier stability beyond bounded targets}
\author{mxym\thanks{Repository-account attribution: \url{https://github.com/mxym/math}. Prepared with AI assistance. No institutional affiliation is asserted.}}
\date{October 7, 2026\quad Version 1}
\hypersetup{pdfauthor={mxym},pdflang={en-US}}
\begin{document}\maketitle
\begin{abstract}
We prove a convex-gradient interpolation inequality with an explicit second-moment tail remainder, for densities of finite integrated coordinate-slice variation. It requires no global bound on the gradients. Under a bounded $p$th gradient moment, $p>2$, the resulting potential-to-gradient exponent is $(p-2)/(3p-2)$, and this interpolation exponent is sharp. Combining the inequality with a radius-independent potential stability estimate yields quantitative Brenier-map stability for unbounded targets with any $p>2$, independently of the ambient dimension as a moment threshold. The application to strongly log-concave sources uses the potential estimate in manuscript 001; compact log-concave sources are also covered by completion of its bounded-target estimate. A tail-infimum bound treats general uniformly integrable second moments, and stretched-exponential target tails give one-third stability with a logarithmic factor. We do not claim that the finite-moment exponent is optimal for transport maps themselves.
\end{abstract}
\paragraph{Status and dependencies.}
The interpolation theorem and its sharp example have complete proofs below, independent of the new transport results in the repository. The Brenier applications explicitly use the radius-independent potential estimate from manuscript 001~\cite{M001}; that proof has not been machine-formalized or independently peer reviewed. The bounded-gradient finite-difference mechanism is inherited from OpenAI family 374~\cite{O374} and 001. The additional argument restricts each coordinate slice to the interval where both derivatives are bounded and charges the complement to the actual gradient tails. Earlier potential-to-gradient approaches and unbounded-target results include Delalande--M\'erigot~\cite{DM}. No priority or journal-tier certification is asserted.

\section{A tail-sensitive interpolation theorem}
Let $D\subset\R^d$ be open and convex. Let $r\ge0$ be a probability density, zero almost everywhere outside $D$. For almost every $z\in e_j^\perp$, assume the coordinate slice $t\mapsto r(z+te_j)$ has a representative of bounded variation. Define
\begin{equation}\label{eq:slice-variation}
 V(r)=\frac12\sum_{j=1}^d\int_{e_j^\perp}
                  \TV_{\R}(r(z+\cdot e_j))\,dz<\infty.
\end{equation}
For a $BV$ density this is half the sum of its directional derivative total variations. For a log-concave density with finite slice functional it equals
\begin{equation}\label{eq:logconcave-functional}
 V(r)=\sum_{j=1}^d\int_{e_j^\perp}
                     \operatorname*{ess\,sup}_t r(z+te_j)\,dz,
\end{equation}
since an integrable unimodal slice has total variation twice its essential maximum. Boundary jumps of a compactly supported density are included.

Let $u,v$ be finite convex functions on $D$ with $u-v\in L^2(r)$ and gradients in $L^2(r)$. Put
\[
 \delta=\inf_{a\in\R}\norm{u-v-a}_{L^2(r)},\qquad
 X=\norm{\nabla u-\nabla v}_{L^2(r)},
\]
\[
 \tau_{u,v}(L)=\int_{\{\norm{\nabla u}>L\}}\norm{\nabla u}^2r
             +\int_{\{\norm{\nabla v}>L\}}\norm{\nabla v}^2r.
\]
All unlabelled integrals are with respect to Lebesgue measure on $D$.

\begin{theorem}[Tail-sensitive gradient interpolation]\label{thm:tail}
For every $h,L>0$,
\begin{equation}\label{eq:main-interpolation}
 \boxed{X^2\le12d\,\delta^2h^{-2}+28V(r)L^2h+4\tau_{u,v}(L).}
\end{equation}
Consequently, with
\[
 A_r=3\,2^{-2/3}(12d)^{1/3}(28V(r))^{2/3},
\]
one has
\begin{equation}\label{eq:infimum}
 \boxed{X^2\le\inf_{L>0}\{A_r\delta^{2/3}L^{4/3}+4\tau_{u,v}(L)\}.}
\end{equation}
When $\delta=0$, both conclusions hold with $X=0$.
\end{theorem}

The statement applies to arbitrary densities satisfying~\eqref{eq:slice-variation}; log-concavity is not needed for this theorem. Convexity of the potentials, rather than convexity of every density level set, is what makes the derivative truncation possible.

\begin{lemma}[One interval]\label{lem:interval}
Let $u,v$ be convex and $L$-Lipschitz on a bounded interval $I$, and let $f=u-v-a$. Then for every $h>0$,
\[
 \int_I|u'-v'|^2\le12h^{-2}\int_I f^2+28L^2h.
\]
The same assertion holds for open or half-open intervals and for restrictions whose endpoints are obtained by limits.
\end{lemma}
\begin{proof}
On $I_h=\{t\in I:t+h\in I\}$ write
$D_hu(t)=(u(t+h)-u(t))/h$. Convexity gives, almost everywhere,
\[
 0\le D_hu-u'\le u'(t+h)-u'(t)\le2L.
\]
Integrating the derivative difference and cancelling translates yields at most $2Lh$, so
$\int_{I_h}(D_hu-u')^2\le4L^2h$; likewise for $v$. Also
$\int_{I_h}|D_hf|^2\le4h^{-2}\int_I f^2$. Using
$u'-v'=D_hf-(D_hu-u')+(D_hv-v')$ gives $12h^{-2}\int_I f^2+24L^2h$ on $I_h$. The remaining interval length is at most $h$, and its squared derivative difference is at most $4L^2$. This proves the estimate, including $h\ge|I|$.
\end{proof}

\begin{proof}[Proof of Theorem~\ref{thm:tail}]
We first explain the one-dimensional level-set counting used below. For a nonnegative integrable $BV$ function $w$ on the line, almost every superlevel set agrees up to null sets with a finite union of bounded intervals. If $n_w(s)$ is the number of these intervals, then
\begin{equation}\label{eq:coarea}
 \int_0^\infty n_w(s)\,ds=\tfrac12\TV(w).
\end{equation}
This is the one-dimensional coarea formula~\cite{FR,Fleming}. It can also be seen directly by using
$|b-c|=\int_0^\infty|\ind_{b>s}-\ind_{c>s}|\,ds$ on nested refining partitions in the definition of variation. Monotone convergence gives the integral of the variation of the superlevel indicator; each bounded component contributes two jumps. Integrability rules out an unbounded component at a positive level.

Fix a coordinate line and a positive level $s$ for its density. On each interval component $I$ of that superlevel set, the derivatives of $u$ and $v$ are monotone where they exist. Their common good set
\[
 J=\{t\in I:|u'(t)|\le L,\ |v'(t)|\le L\}
\]
agrees up to a null set with an interval, possibly empty. On $J$ both restricted functions are $L$-Lipschitz. Lemma~\ref{lem:interval} therefore gives
\[
 \int_J|u'-v'|^2\le12h^{-2}\int_J|u-v-a|^2+28L^2h.
\]
On its complement, at least one derivative has magnitude exceeding $L$, and the pointwise bound
\[
 |u'-v'|^2\le4\bigl(|u'|^2\ind_{|u'|>L}
                       +|v'|^2\ind_{|v'|>L}\bigr)
\]
holds: if just one is large it dominates the magnitude of the other; if both are large use $|b-c|^2\le2(b^2+c^2)$. Thus each original component costs at most one boundary term. This is why clipping the derivatives does not create an uncontrolled number of intervals.

Sum over components, integrate over levels and over transverse lines, and use layer cake, Tonelli, and~\eqref{eq:coarea}. The $j$th coordinate contributes
\[
 \int|\partial_j u-\partial_j v|^2r
 \le12h^{-2}\int|u-v-a|^2r+28L^2hV_j(r)
 +4\int\bigl(|\partial_j u|^2\ind_{|\partial_j u|>L}
              +|\partial_j v|^2\ind_{|\partial_j v|>L}\bigr)r,
\]
where $\sum_j V_j(r)=V(r)$. Convex functions are locally Lipschitz in $D$ and differentiable almost everywhere, so the slice derivatives agree with the corresponding gradient components on a common full-measure set. Endpoints on $\partial D$ do not change the integrals. Summing in $j$, the component-tail sum is bounded by the vector tail $\tau_{u,v}(L)$. Infimize in $a$ to obtain~\eqref{eq:main-interpolation}.

For $\delta>0$, minimize the first two terms in $h$; the identity
$\min_{h>0}(Ah^{-2}+Bh)=3\,2^{-2/3}A^{1/3}B^{2/3}$ yields~\eqref{eq:infimum}. If $\delta=0$, take $h\downarrow0$ in~\eqref{eq:main-interpolation}, then $L\uparrow\infty$. The $L^2$ gradient assumption gives $\tau(L)\to0$, hence $X=0$. The nonnegative density has $V(r)>0$; otherwise all coordinate slices would be constant and integrability would force zero mass.
\end{proof}

\section{Finite moments and a sharp interpolation exponent}
\begin{theorem}\label{thm:moment}
Suppose, in addition, that for some $p>2$
\[
 M=\int(\norm{\nabla u}^p+\norm{\nabla v}^p)r<\infty.
\]
Set $\alpha_p=(p-2)/(3p-2)$. Then
\begin{equation}\label{eq:moment}
 X\le\sqrt2\,A_r^{\frac{3(p-2)}{2(3p-2)}}
                 (4M)^{\frac2{3p-2}}\delta^{\alpha_p}.
\end{equation}
The exponent $\alpha_p$ cannot be replaced by a larger exponent in a uniform inequality of this kind with a fixed source density and bounded $M$.
\end{theorem}
\begin{proof}
The tail is bounded by $\tau(L)\le ML^{2-p}$. For $M\delta>0$, balance the two terms in~\eqref{eq:infimum} by taking
\[
 L=\left(\frac{4M}{A_r\delta^{2/3}}\right)^{1/(p-2/3)}.
\]
Each term then equals
$(A_r\delta^{2/3})^{3(p-2)/(3p-2)}(4M)^{4/(3p-2)}$.
Taking the square root proves~\eqref{eq:moment}; zero cases follow by the previous theorem.

For sharpness take $r=\ind_{(0,1)}$, $v=0$, and
\[
 u_a(x)=a^{-1/p}(x-1+a)_+,\quad0<a<1.
\]
These are restrictions of globally finite convex functions. Direct integration gives
\[
 M=1,\quad X=a^{1/2-1/p},\quad
 \delta^2=a^{3-2/p}\left(\frac13-\frac a4\right).
\]
Thus $X/\delta^\alpha$ diverges for any
$\alpha>(1/2-1/p)/(3/2-1/p)=\alpha_p$. At $p=2$ the same example has $X=1$ while $\delta\to0$, showing that a bounded second gradient moment alone does not supply a uniform potential-to-gradient modulus. This last statement is about arbitrary convex gradients, not about stability in Wasserstein distance.
\end{proof}

\begin{remark}[Sharpness is not transferred automatically to transport]
In the one-dimensional example, one induced target is a point mass and the common-source optimal coupling gives $W_2=X$. It is therefore not a counterexample to a better Brenier-map exponent in $W_2$. We claim sharpness only for the interpolation theorem. A smooth positive density on $[0,1]$, including a truncated Gaussian, gives the same interpolation exponents after normalizing the ramp height by its $p$th moment, since its mass and weighted polynomial integrals near the endpoint have the corresponding leading powers of $a$.
\end{remark}

\section{Application to unbounded-target optimal transport}\label{sec:transport}
For an absolutely continuous probability source $\rho=r\,dx$, write $T_\mu=\nabla U_\mu$ for its quadratic Brenier map and centered convex potential. We separate the transport input from the interpolation theorem:
\begin{equation}\tag{P}\label{eq:P}
 \norm{U_\mu-U_\nu}_{L^2(\rho)}
 \le K_\rho W_2(\mu,\nu)\qquad(\mu,\nu\in\Ptwo(\R^d)).
\end{equation}
The centered representatives are assumed finite on the open convex source domain and in $L^2(\rho)$.

\begin{theorem}[Tail modulus for transport]\label{thm:transport}
Assume~\eqref{eq:slice-variation} and~\eqref{eq:P}. Write $w=W_2(\mu,\nu)$ and
\[
 \tau_{\mu,\nu}(L)=\int_{|y|>L}|y|^2\,d\mu(y)
                    +\int_{|y|>L}|y|^2\,d\nu(y).
\]
Then
\begin{equation}\label{eq:transport-tail}
 \boxed{\norm{T_\mu-T_\nu}_{L^2(\rho)}^2
 \le\inf_{L>0}\{A_rK_\rho^{2/3}w^{2/3}L^{4/3}
                         +4\tau_{\mu,\nu}(L)\}.}
\end{equation}
If $M_p=(\int|y|^p\,d\mu+\int|y|^p\,d\nu)^{1/p}<\infty$ for $p>2$, then
\begin{equation}\label{eq:transport-moment}
 \boxed{\norm{T_\mu-T_\nu}_{L^2(\rho)}
 \le C_{p,r,K_\rho,d}\,M_p^{\frac{2p}{3p-2}}
                         w^{\frac{p-2}{3p-2}}.}
\end{equation}
One may take
$C_{p,r,K_\rho,d}=\sqrt2 A_r^{3(p-2)/(2(3p-2))}4^{2/(3p-2)}K_\rho^{\alpha_p}$.
\end{theorem}
\begin{proof}
The pushforward identity $(T_\mu)_\#\rho=\mu$ identifies gradient tails and moments with target tails and moments. Apply Theorems~\ref{thm:tail} and~\ref{thm:moment}, then~\eqref{eq:P}. No target truncation is used in this step, and no unbounded Lipschitz constant is inserted into a bounded-target estimate.
\end{proof}

\begin{corollary}[Sources supplied by manuscript 001]\label{cor:sources}
The conclusions of Theorem~\ref{thm:transport} hold in either of the following cases:
\begin{enumerate}
\item $r=Z^{-1}e^{-\kappa|x|^2/2-W(x)}$, where $\kappa>0$ and $W$ is proper lower-semicontinuous convex with full-dimensional effective domain. Take $K_\rho=\sqrt{2/\kappa}$.
\item $\rho$ is full-dimensional log-concave, supported on a compact convex set in $B_R$. Take $K_\rho=(1+\sqrt{162})R$.
\end{enumerate}
No common bounded support of the targets, lower atom-mass bound, target density, or moment assumption $p>d$ is required.
\end{corollary}
\begin{proof}
In case (1), manuscript 001 v2, Theorem 1.1, supplies~\eqref{eq:P} for all finite-second-moment targets. Its density-envelope argument gives finite~\eqref{eq:logconcave-functional}. In case (2), its compact companion theorem supplies the same estimate uniformly over bounded targets with a constant independent of the target radius. The following completion argument extends that estimate to $\Ptwo$.

Project a target $\mu$ radially onto balls of radius $n$, obtaining $\mu_n\to\mu$ in $W_2$. The centered convex potentials are Cauchy in $L^2(\rho)$ by the bounded-target estimate. On every compact subset of the interior source domain, the density is positive and bounded below, so they converge in $L^1_{\rm loc}$. The elementary compactness theorem for finite convex functions then gives a finite convex representative of the limit, local uniform convergence, and almost-everywhere gradient convergence. For completeness, local $L^1$ control bounds a convex function above by its ball average; a reflection about a nearby point with bounded value bounds it below on a smaller ball. The resulting local oscillation bound bounds subgradients on a further smaller ball. Arzel\`a--Ascoli and the supporting-plane inequality give the stated conclusions.

Testing gradients against bounded continuous functions identifies their limiting law as $\mu$. The gradient is in $L^2(\rho)$ because that law has finite second moment. Extend the finite convex representative by its supporting affine functions. Fenchel equality on its gradient gives
$U^*(\nabla U)=x\cdot\nabla U-U$, which is integrable in absolute value under the source, using its second moment and $U\in L^2$. Fenchel's inequality is therefore integrable against any competing coupling and proves quadratic optimality of the gradient plan. Almost-everywhere differentiability gives uniqueness. Potentials with the same gradient differ by a constant on the connected source interior. Performing the construction for two targets and passing the bounded-target inequality to the limit establishes~\eqref{eq:P}. Compact log-concavity also gives finite slice functional, as in 001. This completion uses the source-dependent potential estimate; it is not an independent proof of that estimate.
\end{proof}

\begin{corollary}[Uniform integrability and linearized transport]
On any class $\mathcal C\subset\Ptwo(\R^d)$ with
\[
 \lim_{L\to\infty}\sup_{\mu\in\mathcal C}
                  \int_{|y|>L}|y|^2\,d\mu(y)=0,
\]
Theorem~\ref{thm:transport} gives a uniform map-continuity modulus through~\eqref{eq:transport-tail}. Moreover
\[
 W_2(\mu,\nu)\le\norm{T_\mu-T_\nu}_{L^2(\rho)},
\]
so on bounded $p$-moment classes, $p>2$, the common-source linearized transport distance is bi-H\"older equivalent to $W_2$, with upper exponent $\alpha_p$.
\end{corollary}
\begin{proof}
Choose $L$ to control both tails uniformly, and then choose $w$ small to control the first term. For the lower bound, $(T_\mu,T_\nu)_\#\rho$ is a coupling of the two targets. Its cost bounds the optimal coupling cost from above.
\end{proof}

\section{Stretched-exponential tails}
\begin{corollary}\label{cor:exponential}
Under the assumptions of Theorem~\ref{thm:transport}, suppose for some $q,K>0$ and $B<\infty$ that
\[
 \int e^{(|y|/K)^q}\,d\mu(y)+\int e^{(|y|/K)^q}\,d\nu(y)\le B.
\]
For $0<w=W_2(\mu,\nu)\le K$,
\[
 \norm{T_\mu-T_\nu}_{L^2(\rho)}
 \le C_{q,B,r,K_\rho,d}\,K^{2/3}w^{1/3}
                      \bigl(1+\log(K/w)\bigr)^{2/(3q)}.
\]
In particular, a common subgaussian target-tail bound gives a logarithmic factor with exponent $1/3$.
\end{corollary}
\begin{proof}
The finite constant $c_q=\sup_{s\ge0}s^2e^{-s^q/2}$ gives
\[
 \tau_{\mu,\nu}(L)\le c_q B K^2 e^{-(L/K)^q/2}.
\]
Indeed, for $s>L/K$, $s^2\le c_qe^{s^q/2}\le c_qe^{s^q}e^{-(L/K)^q/2}$, and the assumed exponential moment can be integrated. Put $z=w/K\in(0,1]$ and take
\[
 L=K\bigl(2+\tfrac43\log(1/z)\bigr)^{1/q}.
\]
Then the tail is at most $c_qB e^{-1}K^2z^{2/3}$. The first term of~\eqref{eq:transport-tail} is
$A_rK_\rho^{2/3}K^2z^{2/3}(2+\tfrac43\log(1/z))^{4/(3q)}$.
Taking square roots proves the assertion. No claim is made that the logarithmic factor is sharp.
\end{proof}

\section{Scope of verification and of the claimed advance}
The exact verifier accompanying this manuscript checks the boundary-ramp identities, the exponent-balancing identities, and one-dimensional piecewise-affine instances of~\eqref{eq:main-interpolation}, with rational slopes, breakpoints, thresholds, and step sizes. It uses integer and rational arithmetic only. These tests do not formalize the coarea argument or the potential estimate imported from 001.

The sharp statement is the potential-to-gradient interpolation exponent under a $p$th moment bound. The finite-moment \emph{transport} exponent is an established upper estimate in this draft, not an optimality theorem. The second-moment interpolation counterexample must not be presented as a failure of uniform Brenier stability on a bounded $W_2$ set. For comparison, Delalande--M\'erigot~\cite{DM} treats unbounded targets under moment conditions including $p>d$ with a compact source bounded above and below; our corollary changes both the source hypotheses and the moment threshold and uses a different potential estimate. A comprehensive comparison with all later papers has not been completed.

\clearpage
\begin{thebibliography}{9}
\bibitem{M001} mxym, \emph{Sharp one-third stability of Brenier maps for strongly log-concave sources}, v2, October 7, 2026, manuscript 001, \texttt{mxym/math}. Audited source blob: \texttt{d6896f8ac49e3a945324d5d57c54e1797775b1e5}; pinned repository snapshot \texttt{5c6c088aa5abf1c1a4bdca6a8ce5beaa27faaef3}. AI-assisted research draft.
\bibitem{O374} OpenAI, \emph{Sharp One-Third Stability of Brenier Maps}, family 374, \texttt{openai/math}, pinned snapshot \texttt{adc7f1241b42e322a6451854ab7e4b4c146bf78a}. The required finite-difference argument is reproduced above; this item is attributed through manuscript 001.
\bibitem{DM} A. Delalande and Q. M\'erigot, \emph{Quantitative stability of optimal transport maps under variations of the target measure}, Duke Mathematical Journal (2023), DOI: \href{https://doi.org/10.1215/00127094-2022-0106}{10.1215/00127094-2022-0106}; arXiv:2103.05934.
\bibitem{FR} W. H. Fleming and R. Rishel, \emph{An integral formula for total gradient variation}, Archiv der Mathematik 11 (1960), 218--222.
\bibitem{Fleming} W. H. Fleming, \emph{Functions whose partial derivatives are measures}, Illinois Journal of Mathematics 4(3) (1960), 452--478, DOI: \href{https://doi.org/10.1215/ijm/1255456061}{10.1215/ijm/1255456061}.
\end{thebibliography}
\end{document}
